\begin{abstract}
This is a plain-language walkthrough of Vibe Theory, written for one who does not necessarily know all the mathematics,
but still would like to go all the way through to see how the model works.
It builds the whole picture from the ground up: the one idea that
reality is a single growing crystal of experience, the small kit of geometric tools it rests on,
the model itself, and how the physical world and the mind emerge from it, ending with an honest
account of what is solid and what is open. It is deliberately light, a quick tour rather than a
technical treatment, with no equations and no theorems, only the gist of each idea built up so the
full picture can be held at once. The rigorous detail, the eighty-three findings and the geometry,
is in the technical companion \cite{pollard2026techcompanion} and the codebase
\cite{pollard2026vibetest}, and the formal conceptual account is in the conceptual model
\cite{pollard2026vibeconcept}.
\medskip

\noindent\textbf{Status and caveats.} This is exploratory work, a feasibility study rather than a validated result. The aim is to see whether a model of this kind can be made to hang together at all, and so far it looks promising. Much validation and further work remain. Several of the experiments reported here are preliminary, and some may turn out to be inaccurate or improperly set up. Nothing here should be read as established or taken too rigorously at this stage. It is offered in the spirit of curiosity and invitation, a direction that looks worth exploring rather than a finished theory or a claim of discovery.
\end{abstract}
